\chapter{Preparing repositories for AI coding agents}

\section{Engineering principles}

An agent-friendly repository makes engineering constraints discoverable and executable. Clear module boundaries reduce the number of files required to reason about a change. Explicit public interfaces, stable naming, documented architecture and reproducible dependencies reduce accidental coupling. A single deterministic verification entry point reduces ambiguity over what constitutes a completed task.\src{O7}

\section{Implementation approach}

Create a short root instruction file, a component map, a development bootstrap command, a test command and a failure troubleshooting section. Separate application logic from I/O adapters; provide fixture data and avoid live production dependencies. Pin dependencies and keep generated output separate from authored source. Protect CI configuration and acceptance tests from the task identity. Use an architecture test to enforce forbidden imports.

\section{Verification and limitations}

From a fresh checkout, a maintainer should be able to run the same checks without undocumented state. Review changes to build tooling separately; agent instructions cannot substitute for dependency locks or CI enforcement.

\section{Chapter conclusion}

Apply the stated verification criteria before expanding the scope of this workflow. Record deviations, evidence and unresolved risks in the task record.
